\chapter{Conclusiones y Trabajos Futuros}\label{chap:conclusiones}

En este capítulo se presentan las conclusiones obtenidas a partir del desarrollo de la presente tesis, organizadas conforme a los objetivos específicos definidos en el Capítulo~1. Se describen los principales problemas encontrados, se identifican limitaciones del trabajo y se proponen líneas de investigación futura.

\section{Conclusiones por Objetivos Específicos}

\begin{itemize}
    \item \textbf{OE1 — Estado del arte:} Se realizó un análisis crítico de los enfoques actuales en \ac{AQA} y de las técnicas de compresión y destilación aplicadas a video. Se identificó que la literatura previa parte de la suposición implícita de una brecha sustancial entre \textit{Teachers} 3D pesados y \textit{Students} livianos, suposición formulada en un contexto donde los pesos preentrenados, los optimizadores y las prácticas de entrenamiento eran menos maduros.

    \item \textbf{OE2 — Pipeline reproducible:} Se construyó un \textit{pipeline} liviano basado en TSM-MobileNetV2 y MobileNetV3, inicializados con pesos preentrenados de ImageNet y entrenados con prácticas modernas (AdamW, programación cosenoidal, precisión mixta, acumulación de gradientes). El \textit{pipeline} es reproducible (semillas explícitas, configuraciones \texttt{YAML}) y se aplica de forma uniforme a tres dominios.

    \item \textbf{OE3 — Cuantificación de la brecha:} En AQA-7, el \textit{pipeline} liviano alcanza \ac{SRCC} de $0.9021 \pm 0.005$ (TSM-MobileNetV2) y $0.8907 \pm 0.005$ (MobileNetV3) sobre tres semillas, frente a $0.9052$ (\ac{I3D}) y $0.9158$ (SlowFast-R50). La brecha respecto a ambos \textit{Teachers} 3D es menor a $0.025$ \ac{SRCC} y se mantiene robusta entre semillas. En MTL-AQA y JIGSAWS la brecha es similar ($<0.01$ \ac{SRCC}), confirmando que el resultado es estable en dominios de naturaleza distinta.

    \item \textbf{OE4 — Ablaciones de componentes:} El preentrenamiento ImageNet aporta $\sim 0.031$ \ac{SRCC} sobre el backbone aleatorio (E6), constituyendo el componente más importante del \textit{pipeline}. El módulo Temporal Shift (\ac{TSM}) aporta $\sim 0.010$ \ac{SRCC} adicional (E7), beneficio modesto pero positivo. El \textit{pipeline} se sostiene en dos componentes identificables: preentrenamiento ($\sim 3$ puntos) y modelado temporal eficiente ($\sim 1$ punto).

    \item \textbf{OE5 — Eficiencia computacional:} Los \textit{Students} usan menos del $9\%$ de los \ac{FLOPs} del \textit{Teacher} \ac{I3D} (TSM-MobileNetV2: 20.0~G; MobileNetV3: 14.3~G; vs 228.3~G del \ac{I3D}), con parámetros aproximadamente $10\times$ menores y latencia de inferencia entre $2.5$ y $3.4\times$ menor. MobileNetV3 alcanza 39.9~ms por clip de 64 \textit{frames} en RTX 3060 Mobile, habilitando aplicaciones en tiempo real.

    \item \textbf{OE6 — Límites de aplicabilidad:} La transferencia cruzada \textit{zero-shot} obtiene \ac{SRCC} parcial entre dominios deportivos ($\sim 0.55$) y fracasa entre deporte y cirugía. La introducción de \textit{Test-Time Adaptation} por recalibración de BatchNorm (primera aplicación documentada de TTA a \ac{AQA}) mejora 11 de 12 configuraciones evaluadas, con ganancias particularmente notables en pares de mayor \textit{domain shift} (hasta $+0.12$ \ac{SRCC}), sin requerir reentrenamiento ni etiquetas del dominio \textit{target}. El análisis marginal de la destilación muestra que, sobre la línea base liviana fuerte, la \ac{KD} no aporta valor consistente: sólo 1 de 6 configuraciones mejora con \ac{KD}, lo que delimita el alcance del \textit{pipeline} y replantea el rol de la destilación en \ac{AQA} eficiente.
\end{itemize}

\section{Conclusión General}

La presente tesis sostiene tres conclusiones positivas y empíricamente respaldadas:

\textbf{Primero}, un \textit{pipeline} moderno de \textit{Students} livianos (TSM-MobileNetV2 y MobileNetV3 con preentrenamiento ImageNet, AdamW, cosine annealing, AMP y \textit{gradient accumulation}) alcanza \ac{SRCC} a menos de $0.025$ del \textit{Teacher} 3D moderno SlowFast-R50, y a menos de $0.015$ del \ac{I3D}, utilizando el $9\%$ de sus \ac{FLOPs}. Este resultado es estable entre semillas (std $\approx 0.005$) y reproducible en tres dominios (AQA-7, MTL-AQA, JIGSAWS).

\textbf{Segundo}, las ablaciones identifican los componentes responsables del rendimiento del \textit{pipeline}: el preentrenamiento ImageNet (E6, $\sim 3$ puntos de \ac{SRCC}) y el módulo Temporal Shift (E7, $\sim 1$ punto adicional). Esto convierte el \textit{pipeline} en una propuesta interpretable: su aporte no es un efecto agregado opaco, sino la contribución cuantificable de elementos identificables.

\textbf{Tercero}, el análisis marginal de la destilación de conocimiento muestra que, cuando el \textit{Student} se entrena con un \textit{pipeline} fuerte, la \ac{KD} no aporta valor consistente y puede incluso degradar el rendimiento; en contraste, la incorporación de \textit{Test-Time Adaptation} por recalibración de BatchNorm mejora la transferencia cruzada en 11 de 12 configuraciones evaluadas, constituyendo la primera aplicación documentada de TTA a \ac{AQA}. Esto cuestiona la premisa heredada de la literatura previa---según la cual el cierre de brecha entre \textit{Teacher} y \textit{Student} requiere destilación---y replantea la prioridad de futuras propuestas en \ac{AQA} eficiente hacia mecanismos sin gradiente como TTA.

\textbf{Cuarto}, la comparación con el estado del arte sitúa el \textit{pipeline} propuesto por encima de varios métodos representativos del campo en AQA-7 (USDL, GAKD, CoRe, TSA-Net, HGCN) usando aproximadamente el $9\%$ de sus \ac{FLOPs}. En MTL-AQA el \textit{pipeline} queda 7--8 puntos por debajo del SOTA, brecha aceptable dado el ahorro computacional y el contexto de despliegue embebido en el que se ubica la propuesta.

En conjunto, la tesis entrega un \textit{pipeline} liviano reproducible, eficiente, competitivo con métodos del estado del arte en su régimen de costo, y con sus componentes, sus límites de aplicabilidad y un mecanismo complementario de adaptación cross-domain claramente delimitados.

\section{Limitaciones del trabajo}

Para una interpretación honesta del aporte se enumeran a continuación las principales limitaciones, cada una de las cuales abre una línea concreta de trabajo futuro:

\begin{itemize}
    \item \textbf{Réplicas por semilla limitadas a un dataset.} Las réplicas con tres semillas se ejecutaron únicamente sobre AQA-7 (dataset principal). En MTL-AQA y JIGSAWS los resultados se reportan con semilla 42 por costo computacional acumulado. Los hallazgos cuantitativos en esos dos dominios deben leerse considerando esta limitación.

    \item \textbf{Sin validación en hardware embebido real.} La medición de eficiencia se realizó en RTX 3060 Mobile como aproximación a hardware de bajo costo, pero no se ejecutó el modelo en plataformas embebidas reales (Jetson Nano, dispositivos móviles). La validación en tales plataformas, incluyendo consumo energético y memoria de \textit{runtime}, queda como trabajo futuro.

    \item \textbf{Restricciones de memoria que afectaron el régimen de entrenamiento.} El \textit{batch size} físico de 1--2 (impuesto por los 6~GB de \ac{VRAM}) hizo necesario congelar las estadísticas de \textit{BatchNorm} del \textit{Student} durante el entrenamiento con \ac{KD}. Aunque preserva la estabilidad numérica, este procedimiento podría sesgar la evaluación del análisis marginal de \ac{KD} frente a configuraciones con mayor capacidad de hardware.

    \item \textbf{Generalización cross-domain limitada al régimen \textit{zero-shot}.} La transferencia AQA-7 $\to$ JIGSAWS se evaluó sin \textit{fine-tuning}; el resultado negativo no descarta que la destilación o el \textit{pipeline} puedan funcionar tras adaptación de dominio.

    \item \textbf{Subconjunto del paisaje de Students livianos.} Se evaluaron dos arquitecturas (TSM-MobileNetV2 y MobileNetV3); la generalización del hallazgo a otras familias (EfficientNet-Video, X3D-S, MoViNet) requiere validación adicional.
\end{itemize}

\section{Problemas encontrados}

\begin{itemize}
    \item \textbf{Desbalance de clases y distribución de etiquetas:} los datasets empleados presentan concentraciones en rangos medios de puntuación. Esto afectó la capacidad del modelo para distinguir los extremos del espectro de calidad.

    \item \textbf{Variabilidad en las condiciones visuales:} se observaron casos con iluminación deficiente, ángulos de cámara extremos o ruido visual.

    \item \textbf{Limitaciones del hardware de entrenamiento:} la ejecución de modelos 3D (\ac{I3D}, SlowFast) y el cálculo de pérdidas de destilación implicaron una alta demanda de memoria. Esto obligó a reducir tamaños de \textit{batch} y emplear acumulación de gradientes.

    \item \textbf{Ausencia de benchmarks estandarizados para AQA eficiente:} la falta de protocolos unificados dificultó la comparación directa con trabajos previos, haciendo necesario definir un \textit{pipeline} experimental propio.
\end{itemize}

\section{Trabajos futuros}

\begin{itemize}
    \item \textbf{Score Distribution Learning} adaptado al \textit{pipeline} liviano: experimentos preliminares de variantes de regresión distribucional (\textit{Uncertainty-aware Score Distribution Learning}, Tang et al. 2020) mostraron mejoras de hasta $+0.053$ \ac{SRCC} en TSM-MobileNetV2 + JIGSAWS sobre la línea base liviana, sugiriendo que el reformulado del objetivo de regresión es una dirección promisoria.

    \item \textbf{Adaptación de dominio} para resolver el fracaso de transferencia AQA-7 $\to$ JIGSAWS: explorar técnicas como CORAL, MMD o adaptación adversarial.

    \item \textbf{Réplicas por semilla en MTL-AQA y JIGSAWS} para extender la robustez estadística más allá del dataset principal.

    \item \textbf{Implementación en dispositivos reales} (Jetson Nano, móviles) midiendo consumo energético, latencia real y memoria de \textit{runtime}.

    \item \textbf{Aplicación en nuevos dominios} (rehabilitación física, educación física, evaluación de gestos motores en condiciones clínicas).

    \item \textbf{Análisis explicable de decisiones} mediante técnicas de interpretabilidad para facilitar adopción en entornos sensibles como análisis clínico o juicio técnico profesional.
\end{itemize}
